% GENERATED FILE. Edit canonical lessons, then run npm run generate:books.
% canonical-source-hash: fnv1a64:520d55cda63afee4
% canonical-lessons: LA-C324-cunae, LA-C324-puer, LA-C324-infantia, LA-C324-pueritia, LA-C324-educo

\chapter{Puer et pueritia}
\label{ch:la-puer-et-pueritia}
\hlchaptermodality{324}

\begin{chapteropening}
\textbf{By the end of this chapter:} \emph{I can talk about childhood: the cradle, a boy, infancy, childhood and to bring up.}

Complete the last lesson of chapter 324: I can talk about childhood: the cradle, a boy, infancy, childhood and to bring up.
\end{chapteropening}

\section[cūnae, cūnārum]{cūnae, cūnārum — a cradle}

\label{lesson:LA-C324-cunae}



\begin{quote}
Before the new one: say the Latin for to suckle, then the Latin for a nurse.
\end{quote}

\subsection*{You'll want to know: cūnae, cūnārum}
\textbf{cūnae, cūnārum} — \textquotedblleft{}a cradle\textquotedblright{}.

\begin{quote}
Īnfāns in cūnīs dormiēbat.
\end{quote}

\begin{quote}
\emph{The baby was sleeping in the cradle.}
\end{quote}

Feminine, plural. In Plautus a grown man is asked if he still sleeps \emph{in cūnīs}.

\begin{cousinweb}
\textbf{cūnae} is inherited from Proto-Italic \emph{*koinā}, \textquotedblleft{}cradle\textquotedblright{} (EtymDB), from a Proto-Indo-European root meaning \textquotedblleft{}to lie down\textquotedblright{}.
\end{cousinweb}

\begin{grammarlens}[title={What you've built}]
One more word for telling, explaining and describing at B1.
\end{grammarlens}

\subsection*{Your turn}
\begin{itemize}
\raggedright
  \item \emph{Say it:} \emph{cūnae, cūnārum}
  \item \emph{Say it:} \emph{cūnae, cūnārum}, once more
  \item \emph{Say it:} a sentence of your own with \emph{cūnae, cūnārum}
  \item \emph{Recall:} say the Latin for to suckle, then the Latin for a nurse, then say \emph{cūnae, cūnārum} again
\end{itemize}

\begin{tcolorbox}[breakable,colback=teal!4,colframe=teal!35!black,title={Before you move on}]
What does cūnae, cūnārum mean? (\textquotedblleft{}A cradle\textquotedblright{}.) Say it once more.
\end{tcolorbox}
\section[puer, puerī]{puer, puerī — a boy; a child; a slave boy}

\label{lesson:LA-C324-puer}



\begin{quote}
Before the new one: say the Latin for a nurse, then the Latin for a cradle.
\end{quote}

\subsection*{You'll want to know: puer, puerī}
\textbf{puer, puerī} — \textquotedblleft{}a boy; a child; a slave boy\textquotedblright{}.

\begin{quote}
Puer in hortō lūdēbat.
\end{quote}

\begin{quote}
\emph{The boy was playing in the garden.}
\end{quote}

Masculine, 2nd declension; the \emph{e} stays in every form. Also a young slave of any age; the plural \emph{puerī} can mean 'children'.

\begin{cousinweb}
\textbf{puer} is inherited from Proto-Indo-European \emph{*ph\textsubscript{2}weros} (EtymDB), from a root meaning \textquotedblleft{}few, small\textquotedblright{}. English \emph{puerile}, \textquotedblleft{}childish\textquotedblright{}, comes from it.
\end{cousinweb}

\begin{grammarlens}[title={What you've built}]
One more word for telling, explaining and describing at B1.
\end{grammarlens}

\subsection*{Your turn}
\begin{itemize}
\raggedright
  \item \emph{Say it:} \emph{puer, puerī}
  \item \emph{Say it:} \emph{puer, puerī}, once more
  \item \emph{Say it:} a sentence of your own with \emph{puer, puerī}
  \item \emph{Recall:} say the Latin for a nurse, then the Latin for a cradle, then say \emph{puer, puerī} again
\end{itemize}

\begin{tcolorbox}[breakable,colback=teal!4,colframe=teal!35!black,title={Before you move on}]
What does puer, puerī mean? (\textquotedblleft{}A boy; a child; a slave boy\textquotedblright{}.) Say it once more.
\end{tcolorbox}
\section[īnfantia, īnfantiae]{īnfantia, īnfantiae — infancy, early childhood}

\label{lesson:LA-C324-infantia}



\begin{quote}
Before the new one: say the Latin for a cradle, then the Latin for a boy.
\end{quote}

\subsection*{You'll want to know: īnfantia, īnfantiae}
\textbf{īnfantia, īnfantiae} — \textquotedblleft{}infancy, early childhood\textquotedblright{}.

\begin{quote}
Ab īnfantiā in vīllā avī vīxit.
\end{quote}

\begin{quote}
\emph{From early childhood he lived at his grandfather's villa.}
\end{quote}

Feminine, from \emph{īnfāns} (taught), 'not yet speaking'. English \emph{infancy}.

\begin{cousinweb}
\textbf{īnfantia} is built on \emph{īnfāns}, \textquotedblleft{}unable to speak\textquotedblright{} (EtymDB), from \emph{in-}, \textquotedblleft{}not\textquotedblright{}, and \emph{fārī}, \textquotedblleft{}to speak\textquotedblright{}. English \emph{infancy} comes from it.
\end{cousinweb}

\begin{grammarlens}[title={What you've built}]
One more word for telling, explaining and describing at B1.
\end{grammarlens}

\subsection*{Your turn}
\begin{itemize}
\raggedright
  \item \emph{Say it:} \emph{īnfantia, īnfantiae}
  \item \emph{Say it:} \emph{īnfantia, īnfantiae}, once more
  \item \emph{Say it:} a sentence of your own with \emph{īnfantia, īnfantiae}
  \item \emph{Recall:} say the Latin for a cradle, then the Latin for a boy, then say \emph{īnfantia, īnfantiae} again
\end{itemize}

\begin{tcolorbox}[breakable,colback=teal!4,colframe=teal!35!black,title={Before you move on}]
What does īnfantia, īnfantiae mean? (\textquotedblleft{}Infancy, early childhood\textquotedblright{}.) Say it once more.
\end{tcolorbox}
\section[pueritia, pueritiae]{pueritia, pueritiae — childhood, boyhood}

\label{lesson:LA-C324-pueritia}



\begin{quote}
Before the new one: say the Latin for a boy, then the Latin for infancy.
\end{quote}

\subsection*{You'll want to know: pueritia, pueritiae}
\textbf{pueritia, pueritiae} — \textquotedblleft{}childhood, boyhood\textquotedblright{}.

\begin{quote}
In pueritiā multās fābulās audiēbam.
\end{quote}

\begin{quote}
\emph{In my childhood I used to hear many stories.}
\end{quote}

Feminine, from \emph{puer}; up to about the seventeenth year.

\begin{cousinweb}
\textbf{pueritia} is built on \emph{puer}, \textquotedblleft{}boy, child\textquotedblright{}, with the ending \emph{-itia} for a state (EtymDB).
\end{cousinweb}

\begin{grammarlens}[title={What you've built}]
One more word for telling, explaining and describing at B1.
\end{grammarlens}

\subsection*{Your turn}
\begin{itemize}
\raggedright
  \item \emph{Say it:} \emph{pueritia, pueritiae}
  \item \emph{Say it:} \emph{pueritia, pueritiae}, once more
  \item \emph{Say it:} a sentence of your own with \emph{pueritia, pueritiae}
  \item \emph{Recall:} say the Latin for a boy, then the Latin for infancy, then say \emph{pueritia, pueritiae} again
\end{itemize}

\begin{tcolorbox}[breakable,colback=teal!4,colframe=teal!35!black,title={Before you move on}]
What does pueritia, pueritiae mean? (\textquotedblleft{}Childhood, boyhood\textquotedblright{}.) Say it once more.
\end{tcolorbox}
\section[ēducō, ēducāre]{ēducō, ēducāre — to bring up, to rear}

\label{lesson:LA-C324-educo}



\begin{quote}
Before the new one: say the Latin for infancy, then the Latin for childhood.
\end{quote}

\subsection*{You'll want to know: ēducō, ēducāre}
\textbf{ēducō, ēducāre} — \textquotedblleft{}to bring up, to rear\textquotedblright{}.

\begin{quote}
Avia eum ēducāvit, quod parentēs aberant.
\end{quote}

\begin{quote}
\emph{His grandmother brought him up, because his parents were away.}
\end{quote}

\emph{ēducō, ēducāre, ēducāvī, ēducātum}, first conjugation, short \emph{u}. Do not confuse \emph{ēdūcō, ēdūcere}, 'to lead out'. English \emph{educate}.

\begin{cousinweb}
\textbf{ēducō} is built on \emph{ēdūcō}, \textquotedblleft{}to lead out, to draw out\textquotedblright{} (Lewis \& Short), and so on \emph{dūcō}, \textquotedblleft{}to lead\textquotedblright{}. English \emph{educate} comes from it.
\end{cousinweb}

\begin{grammarlens}[title={What you've built}]
One more word for telling, explaining and describing at B1.
\end{grammarlens}

\subsection*{Your turn}
\begin{itemize}
\raggedright
  \item \emph{Say it:} \emph{ēducō, ēducāre}
  \item \emph{Say it:} \emph{ēducō, ēducāre}, once more
  \item \emph{Say it:} a sentence of your own with \emph{ēducō, ēducāre}
  \item \emph{Recall:} say the Latin for infancy, then the Latin for childhood, then say \emph{ēducō, ēducāre} again
\end{itemize}

\begin{tcolorbox}[breakable,colback=teal!4,colframe=teal!35!black,title={Before you move on}]
What does ēducō, ēducāre mean? (\textquotedblleft{}To bring up, to rear\textquotedblright{}.) Say it once more.
\end{tcolorbox}
